\documentclass{article}
\usepackage[utf8]{inputenc}
\usepackage{amsmath}
\usepackage{geometry}
\geometry{margin=2cm}
\begin{document}

\section*{Análise e Resolução Motivacional da Questão 111 — ENEM 2024 — Ciências da Natureza}

\subsection*{Interpretação da Questão, Estratégia e Análise}

Esta questão demanda interpretar a relação entre volumes e massas de gases e substâncias líquidas, buscando determinar a massa de metano necessária para substituir um dado número de mols de etanol na geração de energia, considerando as características fornecidas do biogás.

\textbf{Termos-chave da questão:}
\begin{itemize}
  \item "Substituir": calcular o equivalente em metano para substituir 10 mols de etanol energeticamente.
  \item "Calcular": determinar massa numérica do metano necessária.
\end{itemize}

A estratégia principal envolve:
\begin{enumerate}
  \item Calcular massa e volume do etanol para 10 mols.
  \item Usar a relação dada entre volume de biogás e volume de etanol para encontrar o volume de biogás correspondente.
  \item Extrair o volume correspondente do metano no biogás (70\%).
  \item Converter o volume de metano para mols usando volume molar 22 L/mol.
  \item Calcular a massa de metano a partir de seus mols.
\end{enumerate}

\subsection*{Conteúdo Conceitual e Mini Revisão}

\begin{tabular}{|p{4cm}|p{6cm}|p{6cm}|}
\hline
\textbf{Conceito} & \textbf{Ideia-chave} & \textbf{Aplicação na questão} \\
\hline
Massa molar e número de mols & Relaciona número de mols e massa de substância & Cálculo da massa de etanol e massa de metano a partir do número de mols \\
\hline
Densidade & Relação massa-volume de uma substância & Converter massa do etanol em volume \\
\hline
Volume molar de gases & Volume ocupado por mol de gás nas condições indicadas & Converter volume de metano em quantidade de mols \\
\hline
Proporção e regra de três & Relaciona volumes de biogás e etanol para equivalência & Calcular volume total de biogás que substitui o volume de etanol \\
\hline
\end{tabular}

\medskip

O conhecimento do volume molar é essencial para conversão correta do volume de metano em mols; a densidade permite relacionar massa e volume do etanol, e a regra de três estabelece a relação entre volumes equivalentes de combustível.

\subsection*{Resolução e "Pulo do Gato"}

\begin{enumerate}
  \item Determinar massa de 10 mols de etanol:
  \[
    m_{etanol} = n \times M = 10 \times 46 = 460\ g
  \]
  
  \item Converter em volume pela densidade (0,78 g/mL):
  \[
    V = \frac{460}{0,78} \approx 589,74\ mL = 0,5897\ L
  \]
  
  \item Calcular volume de biogás que substitui 0,5897 L de etanol:
  \[
    V_{biog\acute{a}s} = 0,5897\ L \times \frac{1\ m^{3}}{0,59\ L} \approx 1\ m^{3}
  \]
  
  \item Volume de metano no biogás (70\%):
  \[
    V_{metano} = 0,7 \times 1\ m^{3} = 0,7\ m^{3} = 700\ L
  \]
  
  \item Mols de metano:
  \[
    n = \frac{700}{22} \approx 31,82\ mol
  \]
  
  \item Massa de metano:
  \[
    m = 31,82 \times 16 \approx 509\ g
  \]

\item Arredondamento final no valor mais próximo da alternativa: 510 g.
\end{enumerate}

\textbf{Pulo do gato:} Nunca esqueça de aplicar a porcentagem volumétrica de 70\% do metano no biogás antes de converter para o número de mols e massa – ignorar este dado resulta em superestimação da massa necessária.

\subsection*{Armadilhas e Cuidados}

\begin{itemize}
  \item Confundir volumetria do biogás e porcentagem de metano.
  \item Esquecer de converter unidades adequadamente entre mL, L e m$^{3}$.
  \item Não usar densidade para converter massa do etanol em volume.
  \item Ignorar volume molar no cálculo de quantidade de gás em mols.
  \item Arredondar valores prematuramente antes do cálculo final.
\end{itemize}

\subsection*{Código da Questão}

CN\_QUIMESTQ\_001

\bigskip

\hrule

\bigskip

\textit{Esta resolução segue os padrões editoriais oficiais do ENEM e apresenta o raciocínio de forma clara, passo a passo, ressaltando conceitos fundamentais, métodos e cuidados para o entendimento efetivo da questão.}

\end{document}